% Luc Destombe --- licence CC BY-NC-SA 4.0
% https://creativecommons.org/licenses/by-nc-sa/4.0/deed.fr
% Fichier autonome : compiler avec pdflatex, deux fois (signets, nombre de pages).
\documentclass[a4paper,10pt]{article}
\makeatletter
\newif\iflycee@niveaux
\lycee@niveauxtrue
\newif\iflycee@palatino
\lycee@palatinotrue

\RequirePackage[utf8]{inputenc}
\RequirePackage[T1]{fontenc}
\RequirePackage[french]{babel}
\RequirePackage{etoolbox}

\newcommand{\ShorthandsOff}{\@ifundefined{shorthandoff}{}{\shorthandoff{;:!?}}}
\newcommand{\ShorthandsOn}{\@ifundefined{shorthandon}{}{\shorthandon{;:!?}}}
\ShorthandsOff
\AtBeginDocument{\ShorthandsOn}
\AtBeginEnvironment{tikzpicture}{\ShorthandsOff}

\RequirePackage{geometry}
\RequirePackage{fancyhdr}
\RequirePackage{lastpage}
\RequirePackage{multicol}
\RequirePackage{array}
\RequirePackage{enumitem}
\RequirePackage{xcolor}
\RequirePackage{needspace}

\RequirePackage{amsmath,amssymb}
\RequirePackage{mathpazo}
\DeclareMathSizes{10.95}{10.95}{8}{5.5}
\begingroup\catcode`\_=\active \catcode`\^=\active
\gdef_#1{\sb{\scriptscriptstyle #1}}
\gdef^#1{\sp{\scriptscriptstyle\lycee@moinsepais #1}}
\endgroup
\newcommand\lycee@moins{\mathbin{%
  \setbox\z@\hbox{$\m@th\scriptscriptstyle\lycee@moinsorig$}%
  \dimen@=.55\wd\z@ \advance\dimen@-.5pt
  \hbox to.75\wd\z@{\hss$\m@th\scriptscriptstyle\vcenter{\hbox{%
    \tikz\draw[line width=.5pt,line cap=round](0,0)--(\the\dimen@,0);}}$\hss}}}
\begingroup\lccode`\~=`\- \lowercase{\endgroup\let~}\lycee@moins
\newcommand\lycee@moinsepais{\mathcode`\-="8000 }
\AtBeginDocument{\mathchardef\lycee@moinsorig=\mathcode`\-\relax}
\AtBeginDocument{\catcode`\_=12 \mathcode`\_="8000
  \catcode`\^=12 \mathcode`\^="8000 }
\newcommand\lycee@exposants{%
  \fontdimen13\textfont2=.48\fontdimen6\textfont2
  \fontdimen14\textfont2=.48\fontdimen6\textfont2
  \fontdimen15\textfont2=.48\fontdimen6\textfont2
  \fontdimen13\scriptfont2=.48\fontdimen6\scriptfont2
  \fontdimen14\scriptfont2=.48\fontdimen6\scriptfont2
  \fontdimen15\scriptfont2=.48\fontdimen6\scriptfont2}
\every@math@size\expandafter{\the\every@math@size\lycee@exposants}
\newlength{\retraitcalculs}
\setlength{\retraitcalculs}{2em}

\RequirePackage{tikz}
\usetikzlibrary{arrows.meta,calc,babel,shapes.misc}
\RequirePackage[most]{tcolorbox}
\tcbuselibrary{skins,breakable}

\definecolor{coulDef}{HTML}{1F4E79}
\definecolor{coulProp}{HTML}{7030A0}
\definecolor{coulRem}{HTML}{C55A11}
\definecolor{coulEx}{HTML}{2E7D32}
\definecolor{coulExo}{HTML}{B03A2E}
\definecolor{coulPart}{HTML}{1F4E79}
\definecolor{coulMeth}{HTML}{1B6E4F}
\definecolor{coulCourbe}{HTML}{0B5ED7}
\definecolor{coulCourbeB}{HTML}{C00000}
\definecolor{coulCourbeC}{HTML}{2E7D32}
\definecolor{coulGrille}{HTML}{8C8C8C}
\definecolor{coulRep}{HTML}{1B6E4F}
\definecolor{coulBar}{HTML}{2E7D32}

\newcommand{\R}{\mathbb{R}}
\newcommand{\Rs}{\mathbb{R}^{*}}
\renewcommand{\le}{\leqslant}
\renewcommand{\ge}{\geqslant}
\renewcommand{\approx}{\simeq}
\newcommand{\vect}[1]{\overrightarrow{#1}}
\newcommand{\norme}[1]{\left\lVert #1 \right\rVert}
\newcommand{\intervalle}[2]{\left[#1\,;#2\right]}
\RequirePackage{cancel}
\renewcommand{\CancelColor}{\color{coulExo}}
\newcommand{\fois}[1]{\times{\color{coulExo}#1}}
\newcommand{\barre}[1]{\cancel{#1}}

\tikzset{grille/.style={coulGrille,line width=0.4pt}}
\newcommand{\quadrillage}[4]{\draw[grille,step=1] (#1,#2) grid (#3,#4);}
\tikzset{
  croix/.style={cross out,draw,line width=0.6pt,inner sep=0pt,minimum size=4pt},
  croixdroite/.style={croix,rotate=45,line width=0.8pt},
}
\newcommand{\pt}[3]{\path (#1) node[croix]{} node[#2,font=\small,inner sep=2.5pt]{$#3$};}

\newcommand{\lycee@repere}[4]{%
  \draw[grille,step=1] (#1,#2) grid (#3,#4);
  \draw[-{Stealth[length=2mm]},thick] (#1,0)--({#3+0.4},0) node[below left=-1pt]{$x$};
  \draw[-{Stealth[length=2mm]},thick] (0,#2)--(0,{#4+0.4}) node[below left=-1pt]{$y$};
  \node[below left,font=\scriptsize,inner sep=1.5pt] at (0,0) {$0$};}
\newcommand{\lycee@gradx}[1]{\foreach \k/\l in {#1}{%
  \draw (\k,0.07)--(\k,-0.07) node[below,font=\scriptsize,inner sep=1.5pt]{$\l$};}}
\newcommand{\lycee@grady}[1]{\foreach \k/\l in {#1}{%
  \draw (0.07,\k)--(-0.07,\k) node[left,font=\scriptsize,inner sep=1.5pt]{$\l$};}}
\newcommand{\lycee@intff}[2]{\left[#1\,;#2\right]}
\newcommand{\lycee@intfo}[2]{\left[#1\,;#2\right[}
\newcommand{\lycee@intof}[2]{\left]#1\,;#2\right]}
\newcommand{\lycee@intoo}[2]{\left]#1\,;#2\right[}
\newcommand{\lycee@ens}[1]{\left\{#1\right\}}
\AtBeginDocument{%
  \providecommand{\repere}{\lycee@repere}%
  \providecommand{\gradx}{\lycee@gradx}%
  \providecommand{\grady}{\lycee@grady}%
  \providecommand{\Cf}{\mathcal{C}\sb{\scriptscriptstyle f}}%
  \providecommand{\Cg}{\mathcal{C}\sb{\scriptscriptstyle g}}%
  \providecommand{\intff}{\lycee@intff}%
  \providecommand{\intfo}{\lycee@intfo}%
  \providecommand{\intof}{\lycee@intof}%
  \providecommand{\intoo}{\lycee@intoo}%
  \providecommand{\ens}{\lycee@ens}}

\newlist{questions}{enumerate}{3}
\setlist[questions,1]{label=\arabic*\textdegree),leftmargin=*,itemsep=2pt,topsep=3pt}
\setlist[questions,2]{label=\alph*),leftmargin=*,itemsep=1pt,topsep=2pt}
\setlist[questions,3]{label=\roman*),leftmargin=*,itemsep=1pt,topsep=1pt}
\setlist[itemize]{label=\textbullet,leftmargin=*,itemsep=2pt,topsep=3pt}

\newcommand{\besoinplace}[1]{\par\penalty\@M\begingroup
  \dimen@=#1\relax\dimen@ii=\pagegoal\advance\dimen@ii-\pagetotal
  \ifdim\dimen@>\dimen@ii\ifdim\dimen@ii>\z@\vfil\fi\break\fi\endgroup}

\newcommand{\lycee@pied}{\thepage/\pageref{LastPage}}
\newcommand{\entete}[2]{%
  \pagestyle{fancy}\fancyhf{}%
  \fancyhead[L]{\footnotesize\color{coulPart}\textbf{#1}}%
  \fancyhead[R]{\footnotesize\color{coulPart}\textbf{#2}}%
  \fancyfoot[C]{\footnotesize\lycee@pied}%
  \renewcommand{\headrulewidth}{0.4pt}%
  \renewcommand{\headrule}{\hbox to\headwidth{%
    \color{coulPart!40}\leaders\hrule height \headrulewidth\hfill}}}

\RequirePackage{ccicons}
\newcommand{\lycee@licence}{{\scriptsize\color{gray}%
  \copyright\ Luc Destombe --- \ccbyncsa\ CC BY-NC-SA 4.0}}
\newif\iflycee@licence \lycee@licencetrue
\newcommand{\sanslicence}{\lycee@licencefalse}
\AtBeginDocument{\iflycee@licence\AddToHookNext{shipout/foreground}{%
  \put(0,0){\rlap{\kern\dimexpr1in+\hoffset+\oddsidemargin+\textwidth\relax
    \llap{\raisebox{-\dimexpr1in+\voffset+\topmargin+\headheight+\headsep
      +\textheight+\footskip\relax}{\lycee@licence}}}}}\fi}

\newcommand{\formulesserrees}{%
  \setlength{\abovedisplayskip}{3pt}\setlength{\belowdisplayskip}{3pt}%
  \setlength{\abovedisplayshortskip}{2pt}\setlength{\belowdisplayshortskip}{3pt}}

\setlength{\parindent}{0pt}

\RequirePackage[implicit=false,hidelinks,pdfencoding=auto,psdextra,bookmarksopen,
  bookmarksopenlevel=1,pdfcreator={},pdfproducer={}]{hyperref}
\RequirePackage{bookmark}
\hypersetup{pdfauthor={Luc Destombe},pdfkeywords={CC BY-NC-SA 4.0}}
\newcounter{lycee@signet}
\newcommand{\signet}[3][]{\stepcounter{lycee@signet}%
  \edef\lycee@dest{\if\relax\detokenize{#1}\relax
    signet.\arabic{lycee@signet}\else#1\fi}%
  \hypertarget{\lycee@dest}{}\bookmark[level=#2,dest=\lycee@dest]{#3}}
\pdfstringdefDisableCommands{%
  \def\R{R}\def\vect#1{#1}\def\Cf{Cf}\def\Cg{Cg}\def\fois#1{×#1}%
  \def\barre#1{#1}\def\intervalle#1#2{[#1;#2]}\def\intff#1#2{[#1;#2]}%
  \def\textsuperscript#1{#1}\def\dfrac#1#2{#1/#2}\def\frac#1#2{#1/#2}%
  \def\vec#1{#1⃗}}%
\RequirePackage[official]{eurosym}

\geometry{top=0.8cm,bottom=1.3cm,left=1cm,right=1cm,footskip=0.6cm}

\pagestyle{fancy}
\fancyhf{}
\renewcommand{\headrulewidth}{0pt}
\fancyfoot[C]{\footnotesize Page \thepage{} / \pageref{LastPage}}

\newcommand{\titrefiche}[2]{%
  \begin{center}
    \Large\textbf{#1}\\[2pt]
    \large\textbf{#2}\\[2pt]
  \end{center}
  \vspace{0.10cm}}

\setlength{\columnseprule}{0.4pt}
\setlength{\columnsep}{15pt}
\renewcommand{\columnseprulecolor}{\color{black!45}}
\setlength{\multicolsep}{2pt}

\newcounter{partiefiche}
\renewcommand{\thepartiefiche}{\Roman{partiefiche}}
\newif\ifapresbandeau
\newcommand{\lycee@bandeau}[1]{%
  \par\penalty-600
  \vspace{0.08cm}%
  \noindent\colorbox{coulPart}{%
    \parbox{\dimexpr\linewidth-2\fboxsep\relax}{%
      \color{white}\bfseries\raggedright\strut#1\strut}}%
  \nopagebreak\par\nopagebreak
  \vspace{0.06cm}\nopagebreak
  \apresbandeautrue}
\newcommand{\partiefiche}[1]{%
  \refstepcounter{partiefiche}%
  \lycee@bandeau{\signet{1}{\thepartiefiche{} --- #1}\thepartiefiche{} --- #1}}
\newcommand{\partiefichesansnum}[1]{\lycee@bandeau{\signet{1}{#1}#1}}

\newcounter{exo}
\newlength{\espaceexo}\setlength{\espaceexo}{7pt}
\newlength{\espacetitre}\setlength{\espacetitre}{2pt}
\newcommand{\exo}[1][\relax]{%
  \par
  \ifapresbandeau\apresbandeaufalse\else\penalty-150\addvspace{\espaceexo}\fi
  \refstepcounter{exo}%
  \noindent\signet[exercice-\arabic{exo}]{2}{Exercice \arabic{exo}}%
  \textbf{\color{coulExo}\underline{Exercice \theexo}}%
  \ifx\relax#1\relax\else\textbf{ : #1}\fi
  \par\nopagebreak\vspace{\espacetitre}\nopagebreak
  \ignorespaces}

\setlist[enumerate]{nosep,leftmargin=*,itemsep=2pt,topsep=2pt}
\setlist[itemize]{nosep,leftmargin=*,topsep=2pt}
\newcommand{\choix}[3]{\par\hspace*{1em}%
  \makebox[0.3\linewidth][l]{a) #1}\makebox[0.3\linewidth][l]{b) #2}c) #3}
\newcommand{\deux}[2]{\par\hspace*{0.6em}%
  \makebox[0.48\linewidth][l]{#1}#2}
\newcommand{\trois}[3]{\par\hspace*{0.6em}%
  \makebox[0.32\linewidth][l]{#1}\makebox[0.32\linewidth][l]{#2}#3}

  \renewcommand{\exo}[1][\relax]{%
    \ifapresbandeau\apresbandeaufalse\else\par\penalty-150\fi
    \refstepcounter{exo}%
    \vspace{0.05cm}%
    \noindent\signet[exercice-\arabic{exo}]{2}{Exercice \arabic{exo}}%
    \textbf{\color{coulExo}\underline{Exercice \theexo}}%
    \ifx\relax#1\relax\else\textbf{ : #1}\fi%
    \par\nopagebreak
    \ignorespaces}
  \newcommand{\rappel}[1]{%
    \par\vspace{0.06cm}\nopagebreak
    \noindent\fcolorbox{black!35}{black!5}{%
      \parbox{\dimexpr\linewidth-2\fboxsep-2\fboxrule\relax}{%
        \small\textbf{Rappel.} #1}}%
    \par\vspace{0.06cm}\nopagebreak}
  \setlist[enumerate]{nosep,leftmargin=*}
  \setlist[itemize]{nosep,leftmargin=*}
  \setlength{\multicolsep}{3pt plus 1pt minus 1pt}
  \setlength{\premulticols}{10pt}
  \setlength{\postmulticols}{10pt}
  \newlist{expr}{enumerate}{1}
  \setlist[expr]{label={\Alph*\ $=$},nosep,leftmargin=*,itemsep=0pt}

\renewenvironment{center}{\par\addvspace{3pt}\centering}{\par\addvspace{3pt}}
\formulesserrees
\AtBeginDocument{\formulesserrees}
\clubpenalty=10000
\widowpenalty=10000
\displaywidowpenalty=10000
\brokenpenalty=10000
\setlength{\parskip}{0pt}

\RequirePackage{ragged2e}
\setlength{\RaggedRightRightskip}{0pt plus 1fil}
\AtBeginDocument{\RaggedRight\binoppenalty=10000 \relpenalty=10000 }
\g@addto@macro\@parboxrestore{\RaggedRight}
\makeatother
\geometry{top=0.8cm,bottom=1.3cm,left=1cm,right=1cm,footskip=0.7cm,headheight=12pt}

\newtcblisting{python}{%
  listing only,enhanced,
  colback=black!3,colframe=black!35,boxrule=0.5pt,arc=1pt,
  left=4pt,right=4pt,top=1pt,bottom=1pt,
  listing options={language=Python,basicstyle=\ttfamily\small,
    keywordstyle=\color{coulMeth}\bfseries,columns=fullflexible,
    keepspaces=true,showstringspaces=false,aboveskip=0pt,belowskip=0pt}}
\newcommand{\reperesuite}[4]{%
  \draw[grille,step=1] (-1,#2) grid (#1,#3);
  \draw[-{Stealth[length=2mm]},thick] (-1,0)--({#1+0.4},0)
    node[below left=-1pt,font=\footnotesize]{$n$};
  \draw[-{Stealth[length=2mm]},thick] (0,#2)--(0,{#3+0.4})
    node[below left=-1pt,font=\footnotesize]{$#4$};
  \node[below left,font=\scriptsize,inner sep=1.5pt] at (0,0) {$0$};}

\begin{document}
\titrefiche{1\textsuperscript{re} spécialité --- Suites arithmétiques et géométriques}{Fiche d'exercices du chapitre}

\begin{multicols}{2}

\partiefiche{Suites arithmétiques}

\exo[Reconnaître une suite arithmétique]
Pour chaque début de suite, dire si la suite peut être arithmétique et, si oui, donner
sa raison.
\begin{enumerate}[label=\alph*)]
  \item $7$ ;\; $11$ ;\; $15$ ;\; $19$
  \item $30$ ;\; $24$ ;\; $18$ ;\; $12$ ;\; $6$
  \item $1$ ;\; $2$ ;\; $4$ ;\; $7$ ;\; $11$
  \item $2{,}5$ ;\; $3$ ;\; $3{,}5$ ;\; $4$
\end{enumerate}

\exo[Traduire une situation]
Pour chaque situation, donner le premier terme $u_{0}$, la relation de récurrence et la
raison.
\begin{enumerate}[label=\alph*)]
  \item Un taxi facture $4$~€ de prise en charge, puis $1{,}5$~€ par km ; $u_{n}$ est le
        prix pour $n$~km.
  \item Une bougie de $18$~cm raccourcit de $1{,}2$~cm par heure ; $u_{n}$ est sa hauteur
        au bout de $n$ heures.
  \item Un livret contient $300$~€ ; on y ajoute $25$~€ chaque mois.
\end{enumerate}

\exo[Calculer les premiers termes]
Pour chaque suite arithmétique, calculer les quatre termes qui suivent le premier.
\begin{enumerate}[label=\alph*)]
  \item $u_{0}=4$ et $r=3$ \quad b)~$v_{0}=10$ et $r=-2{,}5$ \quad c)~$w_{1}=-7$ et
        $r=6$
\end{enumerate}

\exo[Trois termes à compléter]
Trouver le nombre $x$ pour que $x-1$ ;\; $2x+1$ ;\; $4x-3$ soient trois termes
consécutifs d'une suite arithmétique. Donner alors ces trois termes et la raison.

\partiefiche{Terme général d'une suite arithmétique}

\exo[Exprimer $u_{n}$ en fonction de $n$]
Pour chaque suite arithmétique, exprimer le terme général en fonction de $n$, puis
calculer le terme demandé.
\begin{enumerate}[label=\alph*)]
  \item $u_{0}=5$ et $r=4$ ; calculer $u_{30}$.
  \item $v_{1}=12$ et $v_{n+1}=v_{n}-3$ ; calculer $v_{40}$.
  \item $w_{5}=20$ et $r=2$ ; calculer $w_{25}$.
\end{enumerate}

\exo[Raison à partir de deux termes]
Pour chaque suite arithmétique, trouver la raison, le premier terme de rang $0$, puis
le terme général.
\begin{enumerate}[label=\alph*)]
  \item $u_{3}=14$ et $u_{8}=39$ \quad b)~$v_{2}=30$ et $v_{10}=6$
  \item[c)] $w_{10}=7$ et $w_{30}=12$
\end{enumerate}

\exo[Est-ce un terme de la suite ?]
La suite $(u_{n})$ est arithmétique, de premier terme $u_{0}=8$ et de raison $3$.
\begin{enumerate}
  \item Exprimer $u_{n}$ en fonction de $n$.
  \item Les nombres $2\,024$ et $2\,025$ sont-ils des termes de la suite ?
  \item Quel est le premier terme supérieur à $1\,000$ ?
\end{enumerate}

\exo[Une cotisation]
La cotisation annuelle d'un club coûte $25$~€ en 2026 et augmente de $1{,}5$~€ par an.
On note $a_{n}$ son prix en $2026+n$.
\begin{enumerate}
  \item Exprimer $a_{n}$ en fonction de $n$ et calculer le prix en 2036.
  \item À partir de quelle année la cotisation dépassera-t-elle $50$~€ ?
\end{enumerate}

\partiefiche{Montrer qu'une suite est (ou n'est pas) arithmétique}

\exo[Avec le terme général]
Les suites suivantes sont-elles arithmétiques ? Si oui, donner la raison et le premier
terme.
\begin{multicols}{2}
\begin{enumerate}[label=\alph*)]
  \item $u_{n}=5n+2$
  \item $v_{n}=n^{2}+1$
  \item $w_{n}=(n+4)^{2}-n^{2}$
  \item $t_{n}=\dfrac{3-2n}{4}$
\end{enumerate}
\end{multicols}

\exo[Avec une relation de récurrence]
Les suites suivantes sont-elles arithmétiques ?
\begin{enumerate}[label=\alph*)]
  \item $u_{0}=2$ et $u_{n+1}=u_{n}+7$
  \item $v_{0}=1$ et $v_{n+1}=2v_{n}+1$
  \item $w_{0}=3$ et $w_{n+1}=w_{n}+n$
  \item $t_{n}=2(n+3)-n$
\end{enumerate}

\exo[Des carrés d'allumettes]
On aligne des carrés faits d'allumettes, côte à côte : un carré demande $4$ allumettes,
deux carrés $7$, trois carrés $10$. On note $a_{n}$ le nombre d'allumettes pour $n$
carrés.
\begin{enumerate}
  \item Justifier que $a_{n+1}=a_{n}+3$ et exprimer $a_{n}$ en fonction de $n$.
  \item Combien faut-il d'allumettes pour $50$ carrés ?
  \item Avec $200$ allumettes, combien de carrés peut-on faire au plus ?
\end{enumerate}

\exo[Somme et produit de suites arithmétiques]
Soit $(u_{n})$ et $(v_{n})$ deux suites arithmétiques de raisons $r$ et $r'$.
\begin{enumerate}
  \item Montrer que la suite $(u_{n}+v_{n})$ est arithmétique ; préciser sa raison.
  \item Avec $u_{n}=n$ et $v_{n}=n$, la suite $(u_{n}\times v_{n})$ est-elle
        arithmétique ?
\end{enumerate}

\partiefiche{Variations : croissance linéaire}

\exo[Étudier et représenter]
La suite $(u_{n})$ est définie par $u_{0}=2$ et $u_{n+1}=u_{n}+1{,}5$.
\begin{enumerate}
  \item Donner la nature de la suite et son sens de variation.
  \item Représenter les termes $u_{0}$ à $u_{4}$.
  \item Donner une équation de la droite qui porte ces points.
\end{enumerate}

\exo[Sens de variation]
Donner le sens de variation de chaque suite arithmétique.
\begin{enumerate}[label=\alph*)]
  \item $u_{n}=3-4n$ \quad b)~$v_{0}=-5$ et $v_{n+1}=v_{n}+0{,}2$ \quad c)~$w_{n}=7$
  \item[d)] $t_{3}=10$ et $t_{8}=-5$
\end{enumerate}

\exo[Lire une suite arithmétique]
Les points représentent les premiers termes d'une suite arithmétique $(u_{n})$.
\begin{center}
\begin{tikzpicture}[x=0.4cm,y=0.4cm]
  \reperesuite{5}{0}{8}{u_{n}}
  \gradx{1,2,3,4} \grady{1,2,3,4,5,6,7}
  \foreach \n in {0,1,2,3}
    \path[coulCourbe] (\n,{1+2*\n}) node[croix]{};
\end{tikzpicture}
\end{center}
\begin{enumerate}
  \item Lire $u_{0}$ et la raison.
  \item Exprimer $u_{n}$ en fonction de $n$ et calculer $u_{20}$.
\end{enumerate}

\exo[À partir de quel rang ?]
Les suites sont définies pour tout entier naturel $n$. Pour chaque suite, résoudre
l'inéquation, puis donner le premier rang $n$ à partir duquel elle est vérifiée.
\begin{enumerate}[label=\alph*)]
  \item $u_{n}=4n-9$ : à partir de quel rang a-t-on $u_{n}>150$ ?
  \item $v_{n}=120-7n$ : à partir de quel rang a-t-on $v_{n}<0$ ?
  \item $w_{n}=2{,}5n+1$ : à partir de quel rang a-t-on $w_{n}\ge 60$ ?
\end{enumerate}

\partiefiche{Somme de termes d'une suite arithmétique}

\exo[Sommes d'entiers]
Calculer les sommes suivantes.
\begin{enumerate}[label=\alph*)]
  \item $A=1+2+\dots+50$ \quad b)~$B=1+2+\dots+150$
  \item[c)] $C=5+10+\dots+200$ \quad d)~$D=101+102+\dots+200$
\end{enumerate}

\exo[Termes consécutifs]
\begin{enumerate}[label=\alph*)]
  \item Soit $u_{n}=4n+1$. Calculer $S=u_{0}+u_{1}+\dots+u_{15}$.
  \item Calculer $T=9+16+23+\dots+128$.
  \item Calculer $U=50+47+44+\dots+2$.
\end{enumerate}

\exo[Somme des nombres impairs]
\begin{enumerate}
  \item Calculer $1+3+5+\dots+99$.
  \item Montrer que $1+3+5+\dots+(2n-1)=n^{2}$.
\end{enumerate}

\exo[Les gradins]
La première rangée d'un amphithéâtre compte $20$ places ; chaque rangée compte $3$
places de plus que la précédente. Il y a $25$ rangées.
\begin{enumerate}
  \item Combien de places la dernière rangée compte-t-elle ?
  \item Combien de places l'amphithéâtre compte-t-il en tout ?
\end{enumerate}

\partiefiche{Problèmes et suites arithmétiques}

\exo[Épargner]
Inès épargne $50$~€ en janvier, puis chaque mois $10$~€ de plus que le mois
précédent. On note $u_{n}$ la somme épargnée le mois $n$ (janvier : $n=0$).
\begin{enumerate}
  \item Calculer $u_{1}$ et $u_{2}$ ; donner la nature de la suite.
  \item Exprimer $u_{n}$ en fonction de $n$.
  \item À partir de quel mois épargne-t-elle plus de $200$~€ ?
  \item Combien aura-t-elle épargné en tout sur les $12$ premiers mois ?
\end{enumerate}

\exo[Deux tarifs de vélos en libre-service]
Tarif A : $2$~€ par trajet. Tarif B : abonnement de $15$~€, puis $0{,}80$~€ par trajet.
On note $a_{n}$ et $b_{n}$ le prix payé pour $n$ trajets.
\begin{enumerate}
  \item Exprimer $a_{n}$ et $b_{n}$ ; quelle est leur nature ?
  \item À partir de combien de trajets le tarif B est-il le plus avantageux ?
\end{enumerate}

\exo[Un entraînement progressif]
Un coureur parcourt $3$~km le premier jour, puis $0{,}5$~km de plus chaque jour.
\begin{enumerate}
  \item Quelle distance parcourt-il le $n$-ième jour ?
  \item Montrer que la distance totale au bout de $n$ jours est
        $\dfrac{n(0{,}5n+5{,}5)}{2}$.
  \item Au bout de combien de jours aura-t-il parcouru $100$~km en tout ? On
        résoudra $n^{2}+11n-400\ge 0$.
\end{enumerate}

\exo[La comète de Halley]
On suppose que la comète de Halley revient tous les $76$ ans ; elle est passée en 1758.
On note $p_{n}$ l'année de son $n$-ième retour après 1758 ($p_{0}=1758$).
\begin{enumerate}
  \item Exprimer $p_{n}$ en fonction de $n$. Retrouver le passage de 1986.
  \item En quelle année repassera-t-elle ?
  \item Selon ce modèle, a-t-elle pu passer en 1066, comme le montre la tapisserie de
        Bayeux ? Que penser du modèle ?
\end{enumerate}

\partiefiche{Suites géométriques}

\exo[Reconnaître une suite géométrique]
Pour chaque début de suite, dire si la suite peut être géométrique et, si oui, donner
sa raison.
\begin{enumerate}[label=\alph*)]
  \item $2$ ;\; $6$ ;\; $18$ ;\; $54$ \qquad b)~$64$ ;\; $32$ ;\; $16$ ;\; $8$
  \item[c)] $3$ ;\; $6$ ;\; $10$ ;\; $15$ \qquad d)~$1\,000$ ;\; $1\,100$ ;\;
        $1\,210$ ;\; $1\,331$
\end{enumerate}

\exo[Traduire une évolution en pourcentage]
Pour chaque situation, donner le premier terme $u_{0}$, la relation de récurrence et la
raison.
\begin{enumerate}[label=\alph*)]
  \item Une ville de $25\,000$ habitants gagne $3\,\%$ d'habitants par an.
  \item Un téléphone de $900$~€ perd $20\,\%$ de sa valeur chaque année.
  \item Une culture de $500$ bactéries double toutes les heures.
  \item Un lac de $2\,000\,000$~m$^{3}$ perd $4\,\%$ de son eau chaque mois.
\end{enumerate}

\exo[Calculer les premiers termes]
Pour chaque suite géométrique, calculer les quatre termes qui suivent le premier.
\begin{enumerate}[label=\alph*)]
  \item $u_{0}=5$ et $q=2$ \quad b)~$v_{0}=1\,000$ et $q=0{,}9$ \quad c)~$w_{1}=4$ et
        $q=1{,}5$
\end{enumerate}

\exo[Trois termes en progression géométrique]
Trouver le nombre $x>0$ pour que $x$ ;\; $x+6$ ;\; $x+18$ soient trois termes
consécutifs d'une suite géométrique. Donner alors ces trois termes et la raison.

\partiefiche{Terme général d'une suite géométrique}

\exo[Exprimer $u_{n}$ en fonction de $n$]
Pour chaque suite géométrique, exprimer le terme général en fonction de $n$, puis
calculer le terme demandé.
\begin{enumerate}[label=\alph*)]
  \item $u_{0}=2$ et $q=3$ ; calculer $u_{6}$.
  \item $v_{1}=6$ et $v_{n+1}=2v_{n}$ ; calculer $v_{8}$.
  \item $w_{3}=40$ et $q=0{,}5$ ; calculer $w_{6}$.
\end{enumerate}

\exo[Raison à partir de deux termes]
Pour chaque suite géométrique de raison positive, trouver la raison, le premier terme
de rang $0$, puis le terme général.
\begin{enumerate}[label=\alph*)]
  \item $u_{1}=6$ et $u_{4}=162$ \quad b)~$v_{2}=400$ et $v_{4}=16$
  \item[c)] $w_{3}=5$ et $w_{6}=40$
\end{enumerate}

\exo[Intérêts simples ou composés]
On place $2\,000$~€. Placement S : on ajoute chaque année $60$~€. Placement C : le
capital augmente de $3\,\%$ par an. On note $s_{n}$ et $c_{n}$ les capitaux au bout de
$n$ années.
\begin{enumerate}
  \item Donner la nature de chaque suite et exprimer $s_{n}$ et $c_{n}$.
  \item Comparer les capitaux au bout de $10$ ans (au centime).
\end{enumerate}

\exo[Plier une feuille]
Une feuille de papier a une épaisseur de $0{,}1$~mm ; chaque pliage en deux double
l'épaisseur. On note $e_{n}$ l'épaisseur, en mm, après $n$ pliages.
\begin{enumerate}
  \item Exprimer $e_{n}$ en fonction de $n$ ; calculer $e_{10}$.
  \item Combien de pliages faudrait-il pour dépasser la hauteur de l'Everest
        ($8\,849$~m) ? la distance Terre-Lune ($384\,400$~km) ?
\end{enumerate}

\partiefiche{Montrer qu'une suite est (ou n'est pas) géométrique}

\exo[Avec le terme général]
Les suites suivantes sont-elles géométriques ? Si oui, donner la raison et le premier
terme.
\begin{multicols}{2}
\begin{enumerate}[label=\alph*)]
  \item $u_{n}=5\times 3^{n}$
  \item $v_{n}=4^{n+1}$
  \item $w_{n}=n^{2}+1$
  \item $t_{n}=\dfrac{7}{2^{n}}$
\end{enumerate}
\end{multicols}

\exo[Avec une relation de récurrence]
Les suites suivantes sont-elles géométriques ?
\begin{enumerate}[label=\alph*)]
  \item $u_{0}=3$ et $u_{n+1}=1{,}2\,u_{n}$
  \item $v_{0}=3$ et $v_{n+1}=v_{n}+1{,}2$
  \item $w_{0}=1$ et $w_{n+1}=2w_{n}+1$
  \item $t_{n}=3^{n}+3^{n+1}$
\end{enumerate}

\exo[Arithmétique, géométrique ou aucune des deux ?]
Pour chaque suite, dire si elle est arithmétique, géométrique ou ni l'une ni l'autre, en
justifiant.
\begin{enumerate}[label=\alph*)]
  \item $u_{n}=4n-1$ \quad b)~$v_{n}=6\times 0{,}5^{n}$ \quad c)~$w_{n}=n^{2}$
  \item[d)] $a_{0}=1$ et $a_{n+1}=a_{n}-2$ \quad e)~$b_{0}=1$ et $b_{n+1}=5b_{n}$
\end{enumerate}

\exo[Produit et somme de suites géométriques]
Soit $(u_{n})$ et $(v_{n})$ deux suites géométriques de raisons $q$ et $q'$.
\begin{enumerate}
  \item Montrer que la suite $(u_{n}\times v_{n})$ est géométrique ; préciser sa raison.
  \item Avec $u_{n}=2^{n}$ et $v_{n}=3^{n}$, la suite $(u_{n}+v_{n})$ est-elle
        géométrique ?
\end{enumerate}

\partiefiche{Variations : croissance exponentielle}

\exo[Sens de variation]
Étudier le sens de variation de chaque suite.
\begin{enumerate}[label=\alph*)]
  \item $u_{n}=4\times 1{,}3^{n}$ \quad b)~$v_{0}=800$ et $v_{n+1}=0{,}6\,v_{n}$
  \item[c)] $w_{n}=-3\times 2^{n}$ \quad d)~$t_{n}=2\times 0{,}95^{n}$
\end{enumerate}

\exo[Deux villes]
En 2026, la ville A compte $20\,000$ habitants et en gagne $500$ par an ; la ville B
compte $15\,000$ habitants et sa population augmente de $3\,\%$ par an. On note $a_{n}$
et $b_{n}$ leurs populations en $2026+n$.
\begin{enumerate}
  \item Exprimer $a_{n}$ et $b_{n}$ en fonction de $n$.
  \item À l'aide de la calculatrice, déterminer l'année où la ville B dépassera la
        ville A.
\end{enumerate}

\exo[Linéaire ou exponentielle ?]
On a représenté les premiers termes de deux suites $(a_{n})$ (\texttimes) et $(g_{n})$
(+), l'une arithmétique, l'autre géométrique.
\begin{center}
\begin{tikzpicture}[x=0.4cm,y=0.4cm]
  \reperesuite{5}{0}{9}{y}
  \gradx{1,2,3,4} \grady{1,2,3,4,5,6,7,8}
  \foreach \n in {0,1,...,4} \path[coulCourbe] (\n,{1+\n}) node[croix]{};
  \foreach \n in {0,1,2,3} \path[coulCourbeB] (\n,{2^\n}) node[croixdroite]{};
\end{tikzpicture}
\end{center}
\begin{enumerate}
  \item Identifier chaque suite, donner sa raison et son terme général.
  \item Comparer $a_{10}$ et $g_{10}$.
\end{enumerate}

\exo[La loi de Moore]
En 1971, un microprocesseur comptait $2\,300$ transistors. Selon la loi de Moore, ce
nombre double tous les deux ans. On note $t_{n}$ le nombre de transistors en
$1971+2n$.
\begin{enumerate}
  \item Exprimer $t_{n}$ en fonction de $n$ et calculer $t_{25}$ ; à quelle année
        correspond-il ?
  \item À l'aide de la calculatrice, déterminer l'année où l'on dépasse un milliard de
        transistors.
  \item Pourquoi parle-t-on de croissance exponentielle ?
\end{enumerate}

\partiefiche{Somme de termes d'une suite géométrique}

\exo[Sommes de puissances]
Calculer les sommes suivantes (la dernière arrondie au millième).
\begin{enumerate}[label=\alph*)]
  \item $A=1+2+2^{2}+\dots+2^{11}$ \quad b)~$B=3+3^{2}+\dots+3^{8}$
  \item[c)] $C=1+0{,}2+0{,}2^{2}+\dots+0{,}2^{5}$
\end{enumerate}

\exo[Somme de termes consécutifs]
\begin{enumerate}[label=\alph*)]
  \item Soit $u_{n}=3\times 2^{n}$. Calculer $u_{0}+u_{1}+\dots+u_{7}$.
  \item Soit $v_{n}=1\,000\times 0{,}5^{n}$. Calculer $v_{0}+v_{1}+\dots+v_{5}$.
\end{enumerate}

\exo[Les grains de riz]
On pose $1$ grain de riz sur la première case d'un échiquier, $2$ sur la deuxième, $4$
sur la troisième, et ainsi de suite en doublant.
\begin{enumerate}
  \item Combien de grains y a-t-il sur la $10$\textsuperscript{e} case ?
  \item Combien de grains y a-t-il en tout sur les $10$ premières cases ? sur les $64$
        cases (sous la forme $2^{p}-1$) ?
\end{enumerate}

\exo[L'escargot]
Un escargot parcourt $1$~m la première heure, puis chaque heure la moitié de la
distance de l'heure précédente.
\begin{enumerate}
  \item Montrer que la distance totale au bout de $n$ heures est
        $S_{n}=2\left(1-0{,}5^{n}\right)$.
  \item Calculer $S_{10}$. L'escargot parcourra-t-il un jour $2$~m ?
  \item Conjecturer la limite de $(S_{n})$.
\end{enumerate}

\partiefiche{Problèmes et suites géométriques}

\exo[Les rebonds d'une balle]
Une balle lâchée de $2$~m remonte à chaque rebond à $70\,\%$ de sa hauteur
précédente. On note $h_{n}$ la hauteur, en m, après le $n$-ième rebond ($h_{0}=2$).
\begin{enumerate}
  \item Donner la nature de $(h_{n})$ et exprimer $h_{n}$ en fonction de $n$.
  \item Calculer la hauteur après $5$ rebonds (au cm).
  \item Écrire une fonction \texttt{seuil()} qui renvoie le premier rebond après lequel
        la hauteur est inférieure à $10$~cm, puis donner ce rebond.
\end{enumerate}

\exo[Doubler un placement]
On place $5\,000$~€ à $2{,}5\,\%$ par an (intérêts composés). On note $c_{n}$ le capital
au bout de $n$ années.
\begin{enumerate}
  \item Exprimer $c_{n}$ en fonction de $n$ et calculer $c_{10}$ (au centime).
  \item Au bout de combien d'années le capital aura-t-il doublé ?
\end{enumerate}

\exo[Téléchargements d'une application]
Une application est téléchargée $2\,000$ fois le premier mois, puis le nombre de
téléchargements augmente de $8\,\%$ par mois. On note $u_{n}$ le nombre de
téléchargements le mois $n$ (premier mois : $n=0$).
\begin{enumerate}
  \item Exprimer $u_{n}$ en fonction de $n$ et calculer $u_{11}$.
  \item Calculer le nombre total de téléchargements sur les $12$ premiers mois.
\end{enumerate}

\exo[Deux offres d'emploi]
Offre A : $24\,000$~€ la première année, puis $1\,000$~€ de plus chaque année. Offre B :
$22\,000$~€ la première année, puis $6\,\%$ de plus chaque année.
\begin{enumerate}
  \item Quelle offre rapporte le plus la première année ? la dixième année ?
  \item Calculer, pour chaque offre, le total gagné sur dix ans. Conclure.
\end{enumerate}

\partiefiche{Suites arithmético-géométriques}

\exo[Avec une suite auxiliaire]
La suite $(u_{n})$ est définie par $u_{0}=100$ et $u_{n+1}=0{,}8\,u_{n}+30$. On pose
$v_{n}=u_{n}-150$.
\begin{enumerate}
  \item Calculer $u_{1}$ et $u_{2}$ ; montrer que $(u_{n})$ n'est ni arithmétique ni
        géométrique.
  \item Montrer que $(v_{n})$ est géométrique ; préciser sa raison et son premier terme.
  \item Exprimer $v_{n}$, puis $u_{n}$, en fonction de $n$.
  \item Étudier le sens de variation de $(u_{n})$, puis conjecturer sa limite.
\end{enumerate}

\exo[Une suite qui s'envole]
La suite $(u_{n})$ est définie par $u_{0}=4$ et $u_{n+1}=3u_{n}-4$. On pose
$v_{n}=u_{n}-2$.
\begin{enumerate}
  \item Montrer que $(v_{n})$ est géométrique ; préciser sa raison et son premier terme.
  \item Exprimer $u_{n}$ en fonction de $n$.
  \item Étudier le sens de variation de $(u_{n})$, puis conjecturer sa limite.
\end{enumerate}

\exo[Les poissons d'un étang]
Un étang compte $2\,000$ poissons. Chaque année, $30\,\%$ des poissons disparaissent et
l'on en introduit $450$. On note $u_{n}$ le nombre de poissons au bout de $n$ années.
\begin{enumerate}
  \item Justifier que $u_{n+1}=0{,}7\,u_{n}+450$.
  \item On pose $v_{n}=u_{n}-1\,500$. Montrer que $(v_{n})$ est géométrique, puis que
        $u_{n}=1\,500+500\times 0{,}7^{n}$.
  \item Étudier le sens de variation de $(u_{n})$ et conjecturer sa limite.
  \item Au bout de combien d'années y aura-t-il moins de $1\,550$ poissons ?
\end{enumerate}

\exo[Trouver la suite auxiliaire]
La suite $(u_{n})$ est définie par $u_{0}=2$ et $u_{n+1}=0{,}6\,u_{n}+4$.
\begin{enumerate}
  \item Résoudre l'équation $\ell=0{,}6\,\ell+4$.
  \item On pose $v_{n}=u_{n}-\ell$. Montrer que $(v_{n})$ est géométrique.
  \item En déduire $u_{n}$ en fonction de $n$, puis la limite de $(u_{n})$.
\end{enumerate}

\partiefichesansnum{Exercices type bac}

\exo[Automatismes (QCM)]
Pour chaque question, une seule réponse est exacte.
\begin{enumerate}
  \item La suite arithmétique de premier terme $u_{0}=7$ et de raison $-2$ vérifie :\par
        a)~$u_{10}=-13$ \quad b)~$u_{10}=-27$ \quad c)~$u_{10}=27$ \quad d)~$u_{10}=-20$
  \item Une hausse de $4\,\%$ par an donne une suite géométrique de raison :\par
        a)~$0{,}04$ \quad b)~$4$ \quad c)~$1{,}4$ \quad d)~$1{,}04$
  \item $1+2+3+\dots+40$ vaut :\par
        a)~$800$ \quad b)~$820$ \quad c)~$1\,600$ \quad d)~$1\,640$
  \item $v_{0}=3$ et $v_{n+1}=2v_{n}$ ; alors $v_{n}$ vaut :\par
        a)~$3+2n$ \quad b)~$2\times 3^{n}$ \quad c)~$3\times 2^{n}$ \quad d)~$6^{n}$
  \item $1+3+3^{2}+\dots+3^{5}$ vaut :\par
        a)~$\dfrac{1-3^{6}}{1-3}$ \quad b)~$\dfrac{1-3^{5}}{1-3}$ \quad c)~$3^{6}$ \quad
        d)~$18$
\end{enumerate}

\exo[Un apiculteur]
\textbf{Partie A.} Un apiculteur possède $12$ ruches en 2026 et en installe $5$ de plus
chaque année. On note $r_{n}$ le nombre de ruches en $2026+n$.
\begin{enumerate}
  \item Quelle est la nature de $(r_{n})$ ? Exprimer $r_{n}$ en fonction de $n$.
  \item En quelle année aura-t-il au moins $60$ ruches ?
\end{enumerate}
\textbf{Partie B.} À cause des sécheresses, la production d'une ruche, $25$~kg de miel en
2026, baisse de $4\,\%$ par an. On note $m_{n}$ la production d'une ruche en $2026+n$.
\begin{enumerate}[resume]
  \item Quelle est la nature de $(m_{n})$ ? Exprimer $m_{n}$ en fonction de $n$.
  \item Calculer la production d'une ruche en 2036 (au dixième de kg).
  \item Compléter la fonction pour qu'elle renvoie le nombre d'années au bout duquel la
        production passe sous $20$~kg, puis donner l'année.
\begin{python}
def seuil():
    m = 25
    n = 0
    while .......... :
        m = ..........
        n = n + 1
    return n
\end{python}
  \item Quelle quantité de miel une ruche produit-elle en tout de 2026 à 2035 ?
\end{enumerate}

\exo[Trottinettes en libre-service]
Une ville dispose de $800$ trottinettes en libre-service en janvier. Chaque mois, $5\,\%$
des trottinettes sont retirées (pannes) et l'opérateur en ajoute $60$. On note $u_{n}$ le
nombre de trottinettes au bout de $n$ mois.
\begin{enumerate}
  \item Calculer $u_{1}$, puis justifier que $u_{n+1}=0{,}95\,u_{n}+60$.
  \item On pose $v_{n}=u_{n}-1\,200$. Montrer que $(v_{n})$ est géométrique ; préciser
        sa raison et son premier terme.
  \item En déduire que $u_{n}=1\,200-400\times 0{,}95^{n}$.
  \item Étudier le sens de variation de $(u_{n})$ et conjecturer sa limite.
  \item À l'aide de la calculatrice ou d'un programme, déterminer au bout de combien de
        mois la ville disposera de plus de $1\,000$ trottinettes.
\end{enumerate}

\end{multicols}

\end{document}
